\bookchapter{第二十一章：移動網絡與信息安全深探}

\chapterintro{吾輩日以手機上網、支付、閒談，甚覺便利。然君曾思之乎？君之聊天記錄、支付密碼、照片……此等信息於網上“裸奔”之時多矣？實言之，衆人未嘗思及此問題，然一旦生變——譬如賬號被盜、銀行卡被刷——屆時追悔已晚。本章當詳論，何以使君之數字生活不至如此“透明”。}

\booksection{一、VPN：網絡中之“隱形斗篷”}

VPN，全稱虛擬專用網絡（Virtual Private Network）。君可想象其爲於君家與所欲往之網站間，築一條專屬且加密之地下隧道。君於隧道中往來，外人不辨君手中何物、欲往何處。

\begin{itemize}
  \item \textbf{尋常上網}：君之手機公共網絡目標網站。中間所經之路由器、交換機，任一環節皆或有人"偷窺"君所發之數據。尤以咖啡廳、機場等公共Wi-Fi爲甚，言之難聽，直如"裸奔"。
  \item \textbf{使用VPN}：君之手機VPN加密隧道VPN服務器目標網站。數據經加密封裝，"偷窺者"所見唯亂碼而已。且目標網站所見者乃VPN服務器之IP地址，非君真實之IP。
\end{itemize}

然VPN究能何爲？簡言之：其一，\textbf{護君於公共Wi-Fi上之隱私}；其二，\textbf{繞過地域限制}（譬如觀看僅於特定國家上線之內容）；其三，\textbf{防君之網絡運營商（ISP）記錄君之瀏覽習慣}。然吾須提醒一句：\textbf{切莫用免費VPN！}免費之背後往往意味着君之數據即其“商品”。吾有一友嘗用某免費VPN，結果手機莫名始彈廣告，後察之，乃該VPN將其瀏覽數據打包售予廣告商也。故擇VPN，或擇付費之大品牌（如NordVPN、ExpressVPN），或乾脆不用。

再者，VPN非萬能。其唯能護“傳輸中”之數據，若目標網站本身不安全、或君之手機已中惡意軟件，VPN亦無能爲力。

\booksection{二、公共Wi-Fi之“甜蜜陷阱”——比君所想更險}

君或以爲，不過連一免費Wi-Fi耳，能有多大風險？吾告君一事：安全研究者曾做實驗，於機場設一名爲“Starbucks\_Free”之Wi-Fi熱點，未及一時辰已有上百人連接。實則此熱點乃研究者自設，凡經此熱點之\textbf{未加密流量}——譬如君所訪之HTTP網頁內容、未啓用TLS之App通信——彼等確能截獲而觀之。

然須注意一關鍵前提：今君日常所用之微信、網銀、支付寶，絕大多數已走HTTPS加密。換言之，攻擊者縱截獲數據包，所見亦唯亂碼，\textbf{無法直接讀出君之聊天內容或支付密碼}。然彼等仍可知君訪問了哪些網站（域名），若某些老舊App或網頁尚用HTTP明文傳輸，則內容將盡露。此等攻擊名曰“中間人攻擊”（Man-in-the-Middle），其理乃攻擊者於君與真正Wi-Fi路由器之間“插足”。更可畏者，多人之手機設爲“自動連接已知Wi-Fi”，攻擊者但需將己之熱點名設爲君曾連過之某Wi-Fi（譬如“ChinaNet”或“CMCC”），君之手機即自動連上。

何以防之？數條實用之策：

\begin{itemize}
  \item \textbf{認準HTTPS}：地址欄開頭爲https://，所多之"s"即Secure之意。訪問任何需輸密碼之網站時，若見http://（無s），立止勿輸！今主流瀏覽器皆於非HTTPS網站前顯"不安全"之警，見此標記當警惕。
  \item \textbf{關閉"自動連接"}：無論iPhone抑或安卓，皆於Wi-Fi設置中關閉"自動加入"或"自動連接"。寧每次手動點之，勿使手機自斷。
  \item \textbf{以手機流量行敏感操作}：轉賬、登錄網銀、輸入身份證號……此等操作老老實實用己之4G/5G流量，勿省那點流量錢。
  \item \textbf{忘記不常用之Wi-Fi}：出差歸來、旅遊歸來，將之前臨時連過之酒店Wi-Fi、機場Wi-Fi皆"忘記"之。如此下次手機不會自動連至同名假熱點。
\end{itemize}

\booksection{三、HTTPS何以護君？——SSL/TLS握手}

方纔提及HTTPS，其究竟何以運作？每次君訪問HTTPS網站，瀏覽器與服務器間將行一次“SSL/TLS握手”，其全過程大約如此：

\begin{enumerate}
  \item 君之網絡運營商
  \item 服務器回復："善，吾自君之清單中擇此加密方案，此乃吾之數字證書，內含吾之公鑰，由某某CA機構簽名認證。"
  \item 瀏覽器驗證證書之真僞：為否由可信之CA簽發？為否過期？域名為否相符？
  \item 驗證既過，雙方以公鑰加密協商出一"會話密鑰"，此後之通信皆以此密鑰加密。
\end{enumerate}

此猶如君往銀行辦事：先遞上身份證（證書），櫃員覈實君之身份，而後爾等約定一暗號，此後之對話唯爾等能解。縱有人於旁偷聽，亦唯聞加密後之亂碼。

然須注意：HTTPS唯保“傳輸過程”加密，不保目標網站爲善類。釣魚網站亦可用HTTPS——彼小鎖圖標唯示君所連乃“某個網站”，不代表其即君所欲往者。故仍須自核網址為否無誤。

\booksection{四、DNS安全：君之“網絡電話本”安全乎？}

DNS（域名系統）猶如互聯網之電話本：君輸入“baidu.com”，DNS爲君查得其對應之IP地址（譬如110.242.68.66），而後瀏覽器方能連之。問題於於，傳統DNS查詢乃\textbf{明文}、\textbf{未加密}者！為謂君之ISP、Wi-Fi提供者、或任何中間節點，皆能見君訪問過哪些網站。

更可惡者乃“DNS劫持”：攻擊者於DNS環節做手腳，將“baidu.com”指向一假釣魚網站。君明明輸入正確網址，卻被引至一模一樣然實則虛假之登錄頁面，一輸賬號密碼即遭殃。

幸今已有加密DNS之策：

\begin{itemize}
  \item \textbf{DNS over HTTPS（DoH）}：將DNS查詢藏於HTTPS流量中，混於正常網頁瀏覽流量之內，根本無從分辨何爲DNS查詢。
  \item \textbf{DNS over TLS（DoT）}：用TLS加密協議專門保護DNS查詢，端口獨立，網管能知君用了DoT然觀不到內容。
\end{itemize}

吾私薦用Cloudflare之1.1.1.1或Google之8.8.8.8作DNS服務器，二者皆支持加密DNS。於安卓之“私人DNS”設置中徑填“one.one.one.one”或“dns.google”即可，甚簡。iPhone稍繁，需裝一描述文件或經App設置，然亦值得爲之。

\booksection{五、“無痕瀏覽”真能隱身乎？}

先言結論：\textbf{無痕模式乃君家之窗簾，非隱身斗篷。}

其所作：不留痕跡於本地——關閉無痕窗口後，瀏覽器不存君之瀏覽歷史、Cookies、表單填寫記錄、搜索記錄。借他人電腦查一物、或買禮物不欲爲家人自瀏覽記錄所覺，用無痕模式爲宜。

其不能作：

\begin{itemize}
  \item 君所訪之網站\textbf{照舊能見君來過}，君之IP地址、設備信息、瀏覽行爲，網站端一清二楚。
  \item 君之網絡運營商\textbf{照舊能見君訪問了哪些網站}，唯不見具體內容（若用了HTTPS）。
  \item 公司/學校之網絡管理員\textbf{一樣能監控}君之上網行爲。
  \item 君所下載之文件、所添之書籤將\textbf{保留下來}，不會隨着關閉窗口而消失。
\end{itemize}

故當記住：無痕模式乃“本地擦除器”，非“網絡隱身衣”。若君真需隱藏上網痕跡，須用VPN+Tor瀏覽器之組合拳，然於日常使用而言，知此足矣。

\booksection{六、瀏覽器隱私：勿令瀏覽器賣君}

君或未察，瀏覽器本身即一巨大之信息收集器。Cookies、瀏覽歷史、緩存、自動填充、位置權限……每一樣皆或泄露君之隱私。數條建議：

\begin{itemize}
  \item \textbf{定期清理瀏覽數據}：每隔一段時日，往設置中將Cookies、緩存、歷史記錄清之。非爲防"被追蹤"——彼靠清Cookies已不太管用——乃爲防手機愈來愈卡。瀏覽器緩存可佔數G之空間。
  \item \textbf{勿信"請勿追蹤"之幻想}：瀏覽器中之"發送'請勿追蹤'請求"選項，本質乃"請"網站勿追蹤君。然網站聽否全然另一事，多數網站徑直略之。勿指望之。
  \item \textbf{善用隱私保護瀏覽器}：Firefox Focus、DuckDuckGo瀏覽器、Brave等瀏覽器默認攔截追蹤器與廣告，比Chrome與Safari之默認設置乾淨多矣。
  \item \textbf{搜索引擎亦要緊}：以DuckDuckGo代百度/Google搜索，其不錄君之搜索歷史，不建用戶畫像。實話言，搜索結果確不如Google精準，然日常查資料足矣。
  \item \textbf{檢查App權限}：往手機設置中觀君之瀏覽器被授予何權限。其真需訪問君之位置、麥克風、通訊錄乎？大概率不需。
\end{itemize}

\booksection{七、實用安全清單：自今始護己}

言之多矣，終予君一簡明之安全清單，費半小時了之，君之手機安全性將有質之飛躍：

\begin{enumerate}
  \item \textbf{關閉Wi-Fi自動連接}：iPhone於設置-無線局域網中，安卓於Wi-Fi設置中。
  \item \textbf{換一加密DNS}：安卓徑設私人DNS爲"dns.google"，iPhone以App或描述文件設1.1.1.1之DoH。
  \item \textbf{開啓雙因素認證（2FA）}：微信、支付寶、Apple ID、Google賬號，能開2FA者皆開。短信驗證碼已為最低級別之2FA，能用TOTP（如Google Authenticator、Microsoft Authenticator）更佳。
  \item \textbf{查一遍App權限}：將每個App之權限過一遍，尤以位置、通訊錄、短信、麥克風、相機此五項。君將發現有些手電筒App竟要君之通訊錄——速卸之。
  \item \textbf{定期清理瀏覽器數據}：設一日曆提醒，每月一次，清緩存、Cookies。
  \item \textbf{勿用公共充電樁之數據線}：USB口既能充電亦能傳數據，惡意充電樁（名曰"Juice Jacking"）可趁君充電時竊數據。出門自帶己之充電頭與線。
\end{enumerate}

\begin{tcbpractice}{即刻爲之：爲君之手機做一次"安全體檢"}
\begin{enumerate}
  \item \textbf{DNS升級}：安卓：設置 > 連接與共享 > 私人DNS > 填入"one.one.one.one"或"dns.google"。iPhone：以Safari搜"1.1.1.1 DNS profile"安裝描述文件，或下載Cloudflare之1.1.1.1 App。
  \item \textbf{Wi-Fi審計}：打開Wi-Fi列表，將不識、不常用、名怪之熱點皆"忘記"之。
  \item \textbf{權限大掃除}：入設置 > 隱私/權限管理，逐個App審查權限。標記出權限要求不合理（譬如計算器要通訊錄）之App，考慮卸載或用替代品。
  \item \textbf{檢查賬號安全}：打開Apple ID或Google賬號之安全頁面，觀有無不明設備登錄，確認2FA已開啓，檢查恢復聯繫人為否可信。
\end{enumerate}

\end{tcbpractice}
